% GENERATED FILE. Edit canonical lessons, then run npm run generate:books.
% canonical-source-hash: fnv1a64:5b2dc3c9fef52d04
% canonical-lessons: MR-C413-nivas, MR-C413-dharmasala, MR-C413-visramgrh, MR-C413-mukkam, MR-C413-khanaval

\chapter{A Place to Stay}
\label{ch:mr-a-place-to-stay}
\hlchaptermodality{413}

\begin{chapteropening}
\textbf{By the end of this chapter:} \emph{I can talk about lodging, rest houses, guest houses, a night halt and an eating house.}

Complete the last lesson of chapter 413: I can talk about lodging, rest houses, guest houses, a night halt and an eating house.
\end{chapteropening}

\section[\texorpdfstring{\mr{निवास} (nivās)}{निवास (nivās)}]{\mr{निवास} (nivās) — lodging; a residence}

\label{lesson:MR-C413-nivas}



\begin{quote}
Before the new one: say the Marathi for to stray, then the Marathi for to hang about.
\end{quote}

\subsection*{You'll want to know: \mr{निवास}}
\textbf{\mr{निवास}} — \emph{nivās} — \textquotedblleft{}lodging; a residence\textquotedblright{}.

\begin{quote}
\mr{सहलीत} \mr{आमचा} \mr{निवास} \mr{एका} \mr{जुन्या} \mr{वाड्यात} \mr{होता}.
\end{quote}

\begin{quote}
\emph{sahalīt āmcā nivās ekā junyā vāḍyāt hotā.}
\end{quote}

\begin{quote}
On the trip we stayed in an old mansion.
\end{quote}

Masculine.

A formal word on signs and forms: \mr{निवास} \mr{आणि} \mr{जेवण}, board and lodging.

\begin{cousinweb}
\textbf{\mr{निवास}} (\emph{nivās}) is a tatsama, taken from Sanskrit \emph{nivāsa}, \textquotedblleft{}dwelling, abode\textquotedblright{}, from \emph{ni-vas}, \textquotedblleft{}to dwell\textquotedblright{} (Monier-Williams).
\end{cousinweb}

\begin{grammarlens}[title={What you've built}]
One more word for getting around.
\end{grammarlens}

\subsection*{Your turn}
\begin{itemize}
\raggedright
  \item \emph{Say it:} \emph{nivās}
  \item \emph{Say it:} \emph{nivās}, once more
  \item \emph{Say it:} \emph{sahalīt āmcā nivās ekā junyā vāḍyāt hotā.} — On the trip we stayed in an old mansion.
  \item \emph{Recall:} say the Marathi for to stray, then the Marathi for to hang about, then say \emph{nivās} again
\end{itemize}

\begin{tcolorbox}[breakable,colback=teal!4,colframe=teal!35!black,title={Before you move on}]
What does \mr{निवास} mean? (\textquotedblleft{}Lodging; a residence\textquotedblright{}.) Say it once more.
\end{tcolorbox}
\section[\texorpdfstring{\mr{धर्मशाळा} (dharmaśāḷā)}{धर्मशाळा (dharmaśāḷā)}]{\mr{धर्मशाळा} (dharmaśāḷā) — a rest house; a pilgrims’ lodge}

\label{lesson:MR-C413-dharmasala}



\begin{quote}
Before the new one: say the Marathi for to hang about, then the Marathi for lodging.
\end{quote}

\subsection*{You'll want to know: \mr{धर्मशाळा}}
\textbf{\mr{धर्मशाळा}} — \emph{dharmaśāḷā} — \textquotedblleft{}a rest house; a pilgrims’ lodge\textquotedblright{}.

\begin{quote}
\mr{मंदिराजवळ} \mr{एक} \mr{धर्मशाळा} \mr{आहे}.
\end{quote}

\begin{quote}
\emph{mandirājavaḷ ek dharmaśāḷā āhe.}
\end{quote}

\begin{quote}
There is a rest house near the temple.
\end{quote}

Feminine.

Often near a temple; a bed there is cheap or free.

\begin{cousinweb}
\textbf{\mr{धर्मशाळा}} (\emph{dharmaśāḷā}) is a tatsama, taken from Sanskrit \emph{dharmaśālā}, built from \emph{dharma}, \textquotedblleft{}duty, charity\textquotedblright{}, and \emph{śālā}, \textquotedblleft{}a hall\textquotedblright{} (Monier-Williams): a hall kept open as a good deed.
\end{cousinweb}

\begin{grammarlens}[title={What you've built}]
One more word for getting around.
\end{grammarlens}

\subsection*{Your turn}
\begin{itemize}
\raggedright
  \item \emph{Say it:} \emph{dharmaśāḷā}
  \item \emph{Say it:} \emph{dharmaśāḷā}, once more
  \item \emph{Say it:} \emph{mandirājavaḷ ek dharmaśāḷā āhe.} — There is a rest house near the temple.
  \item \emph{Recall:} say the Marathi for to hang about, then the Marathi for lodging, then say \emph{dharmaśāḷā} again
\end{itemize}

\begin{tcolorbox}[breakable,colback=teal!4,colframe=teal!35!black,title={Before you move on}]
What does \mr{धर्मशाळा} mean? (\textquotedblleft{}A rest house; a pilgrims’ lodge\textquotedblright{}.) Say it once more.
\end{tcolorbox}
\section[\texorpdfstring{\mr{विश्रामगृह} (viśrāmgṛh)}{विश्रामगृह (viśrāmgṛh)}]{\mr{विश्रामगृह} (viśrāmgṛh) — a guest house; a rest house}

\label{lesson:MR-C413-visramgrh}



\begin{quote}
Before the new one: say the Marathi for lodging, then the Marathi for a rest house.
\end{quote}

\subsection*{You'll want to know: \mr{विश्रामगृह}}
\textbf{\mr{विश्रामगृह}} — \emph{viśrāmgṛh} — \textquotedblleft{}a guest house; a rest house\textquotedblright{}.

\begin{quote}
\mr{घाटावर} \mr{सरकारी} \mr{विश्रामगृह} \mr{आहे}.
\end{quote}

\begin{quote}
\emph{ghāṭāvar sarkārī viśrāmgṛh āhe.}
\end{quote}

\begin{quote}
There is a government guest house on the ghat.
\end{quote}

Neuter.

Run by the government for travelling officers and visitors.

\begin{cousinweb}
\textbf{\mr{विश्रामगृह}} (\emph{viśrāmgṛh}) is built from Sanskrit \emph{viśrāma}, \textquotedblleft{}rest\textquotedblright{}, and \emph{gṛha}, \textquotedblleft{}a house\textquotedblright{} (Monier-Williams).
\end{cousinweb}

\begin{grammarlens}[title={What you've built}]
One more word for getting around.
\end{grammarlens}

\subsection*{Your turn}
\begin{itemize}
\raggedright
  \item \emph{Say it:} \emph{viśrāmgṛh}
  \item \emph{Say it:} \emph{viśrāmgṛh}, once more
  \item \emph{Say it:} \emph{ghāṭāvar sarkārī viśrāmgṛh āhe.} — There is a government guest house on the ghat.
  \item \emph{Recall:} say the Marathi for lodging, then the Marathi for a rest house, then say \emph{viśrāmgṛh} again
\end{itemize}

\begin{tcolorbox}[breakable,colback=teal!4,colframe=teal!35!black,title={Before you move on}]
What does \mr{विश्रामगृह} mean? (\textquotedblleft{}A guest house; a rest house\textquotedblright{}.) Say it once more.
\end{tcolorbox}
\section[\texorpdfstring{\mr{मुक्काम} (mukkām)}{मुक्काम (mukkām)}]{\mr{मुक्काम} (mukkām) — an overnight stay; a halt}

\label{lesson:MR-C413-mukkam}



\begin{quote}
Before the new one: say the Marathi for a rest house, then the Marathi for a guest house.
\end{quote}

\subsection*{You'll want to know: \mr{मुक्काम}}
\textbf{\mr{मुक्काम}} — \emph{mukkām} — \textquotedblleft{}an overnight stay; a halt\textquotedblright{}.

\begin{quote}
\mr{आज} \mr{रात्री} \mr{आमचा} \mr{मुक्काम} \mr{एका} \mr{लहान} \mr{गावात} \mr{आहे}.
\end{quote}

\begin{quote}
\emph{āj rātrī āmcā mukkām ekā lahān gāvāt āhe.}
\end{quote}

\begin{quote}
Tonight we halt in a small village.
\end{quote}

Masculine.

\mr{मुक्काम} \mr{करणे}, to stop for the night; \mr{मुक्कामी}, staying over.

\begin{cousinweb}
\textbf{\mr{मुक्काम}} (\emph{mukkām}) is a loan through Persian, usually traced to Arabic \emph{muqām}, \textquotedblleft{}a place where one stays\textquotedblright{}.
\end{cousinweb}

\begin{grammarlens}[title={What you've built}]
One more word for getting around.
\end{grammarlens}

\subsection*{Your turn}
\begin{itemize}
\raggedright
  \item \emph{Say it:} \emph{mukkām}
  \item \emph{Say it:} \emph{mukkām}, once more
  \item \emph{Say it:} \emph{āj rātrī āmcā mukkām ekā lahān gāvāt āhe.} — Tonight we halt in a small village.
  \item \emph{Recall:} say the Marathi for a rest house, then the Marathi for a guest house, then say \emph{mukkām} again
\end{itemize}

\begin{tcolorbox}[breakable,colback=teal!4,colframe=teal!35!black,title={Before you move on}]
What does \mr{मुक्काम} mean? (\textquotedblleft{}An overnight stay; a halt\textquotedblright{}.) Say it once more.
\end{tcolorbox}
\section[\texorpdfstring{\mr{खानावळ} (khānāvaḷ)}{खानावळ (khānāvaḷ)}]{\mr{खानावळ} (khānāvaḷ) — an eating house}

\label{lesson:MR-C413-khanaval}



\begin{quote}
Before the new one: say the Marathi for a guest house, then the Marathi for an overnight stay.
\end{quote}

\subsection*{You'll want to know: \mr{खानावळ}}
\textbf{\mr{खानावळ}} — \emph{khānāvaḷ} — \textquotedblleft{}an eating house\textquotedblright{}.

\begin{quote}
\mr{आम्ही} \mr{खानावळीत} \mr{गरम} \mr{जेवण} \mr{घेतलं}.
\end{quote}

\begin{quote}
\emph{āmhī khānāvaḷīt garam jevaṇ ghetlaṁ.}
\end{quote}

\begin{quote}
We had a hot meal at the eating house.
\end{quote}

Feminine.

A simple place serving home-style meals for a fixed price; before an ending \mr{खानावळी}-.

\begin{cousinweb}
\textbf{\mr{खानावळ}} (\emph{khānāvaḷ}) is a Marathi word for a cheap, homely place to eat; its origin is not settled.
\end{cousinweb}

\begin{grammarlens}[title={What you've built}]
One more word for getting around.
\end{grammarlens}

\subsection*{Your turn}
\begin{itemize}
\raggedright
  \item \emph{Say it:} \emph{khānāvaḷ}
  \item \emph{Say it:} \emph{khānāvaḷ}, once more
  \item \emph{Say it:} \emph{āmhī khānāvaḷīt garam jevaṇ ghetlaṁ.} — We had a hot meal at the eating house.
  \item \emph{Recall:} say the Marathi for a guest house, then the Marathi for an overnight stay, then say \emph{khānāvaḷ} again
\end{itemize}

\begin{tcolorbox}[breakable,colback=teal!4,colframe=teal!35!black,title={Before you move on}]
What does \mr{खानावळ} mean? (\textquotedblleft{}An eating house\textquotedblright{}.) Say it once more.
\end{tcolorbox}
